\documentclass[10pt,a4paper]{article}
\usepackage[a4paper,margin=1.65cm]{geometry}
\usepackage[T1]{fontenc}
\usepackage[utf8]{inputenc}
\usepackage{lmodern}
\usepackage{microtype}
\usepackage{enumitem}
\usepackage{titlesec}
\usepackage{hyperref}
\usepackage{xcolor}
\usepackage{tabularx}
\definecolor{heading}{HTML}{000000}
\definecolor{secondary}{HTML}{000000}
\hypersetup{colorlinks=true,urlcolor=heading,linkcolor=heading}
\pagestyle{empty}
\setlength{\parindent}{0pt}
\setlength{\parskip}{3pt}
\setlist[itemize]{leftmargin=1.3em,itemsep=2pt,topsep=3pt,parsep=0pt}
\titleformat{\section}{\large\bfseries\color{heading}}{}{0pt}{}[\vspace{-3pt}\titlerule]
\titlespacing*{\section}{0pt}{12pt}{6pt}
\newcommand{\job}[4]{\textbf{#1} \hfill \textit{#2}\\\textit{\color{secondary}#3}\hfill\textit{\color{secondary}#4}\vspace{-2pt}}
\newcommand{\skill}[2]{\textbf{#1:} #2\\}
\begin{document}
\begin{center}
{\LARGE\bfseries UDAYAN PRABIR SINHA}\\[4pt]
{\large\color{heading} Senior Embedded Linux \& Platform Software Engineer}\\[7pt]
Sweden \quad | \quad +46-72-594-8719 \quad | \quad \href{mailto:udayanprabirsinha@gmail.com}{udayanprabirsinha@gmail.com}\\[3pt]
\href{https://www.linkedin.com/in/udayan-prabir-sinha-27799580/}{linkedin.com/in/udayan-prabir-sinha-27799580}
\end{center}
\section{Professional Summary}
Senior Embedded Linux and Platform Software Engineer with approximately 10 years of experience spanning Linux kernel development, device drivers, low-level system software, platform bring-up, and pre-silicon hardware emulation for embedded computing platforms. Experienced in system analysis \& troubleshooting, integrating software using Yocto-based workflows, evaluating new kernel technologies, and designing Embedded Linux platform solutions. Currently combining hands-on kernel engineering with technical leadership and architectural responsibilities across cross-functional engineering teams.
\section{Technical Expertise}
\skill{Embedded Linux}{Linux Kernel Development, New Kernel Feature Analysis, Device Drivers, Platform Bring-Up \& Trouble-Shooting}
\skill{System Design \& Modeling}{Containerization Technologies (runtime, orchestration, image build), System Capacity Dimensioning \& Profiling, System Telemetry, SystemC TLM Based Emulations}
\skill{Debugging, Tracing \& Diagnostics}{ftrace/tracefs, procfs, sysfs, perf, LTTng, lockdep, pstore, kdump, dmesg, strace, gdb, Valgrind, Sanitizer, QEMU}
\skill{Programming Languages}{C, C++, Python, Shell Scripting (BASH), Assembly}
\skill{Build Systems}{Yocto Project, Make, Autotools, CMake}
\skill{Development Tools}{Git, JIRA, AI-Enhanced Workflows}
\section{Professional Experience}
\textbf{Ericsson} \hfill Stockholm, Sweden\\
\job{Technical Lead (Product Guardian)}{July 2024 -- Present}{Technical Leadership / Platform Software Architecture}{Concurrent Role}
\begin{itemize}
\item \textbf{Containerization Technologies:} Lead the architecture and development of in-house container runtime, orchestration and image build solutions for embedded Linux platforms, addressing platform-specific performance and security requirements. Also resulted in architectural improvements to software modularity with the aim to reduce distribution upgrade costs.
\item \textbf{Adaptive CPU Dimensioning:} Develop a scalable, adaptive CPU dimensioning architecture using policy-driven resource allocation to balance workload fairness and performance requirements.
\item \textbf{SoC Utilization \& Telemetry:} Architect a SoC utilization monitoring architecture that integrates hardware telemetry, including CPU and cache utilization, resource pressure, and temperature monitoring, into the existing Ericsson product telemetry framework.
\item \textbf{Cross-Functional Leadership:} Provide technical steering, architectural mentorship and project alignment across multiple cross-functional engineering teams.
\end{itemize}
\vspace{5pt}
\job{Software Developer}{October 2021 -- Present}{Linux Kernel \& Device Drivers}{Concurrent Role}
\begin{itemize}
\item \textbf{Kernel Feature Evaluation \& Profiling:} Conduct deep feasibility studies for back-porting and integrating new kernel capabilities (e.g. CPU Isolation, Linux Audit, kdump) on embedded computing platforms. Also includes system profiling studies such as kernel behavior analysis on massively parallel (100+ CPUs) embedded systems.
\item \textbf{Kernel-Space Device Drivers \& Features:} Develop and maintain kernel-space device drivers (e.g. for PCIe endpoints), kernel panic post-mortem diagnostics infrastructure and custom network layers.
\item \textbf{User-Space Device Drivers:} Developed low-level userspace drivers for proprietary SoCs handling early platform initialization (clock and power rail setup), alongside monitoring functionality for SoC bring-up diagnostics.
\textbf{Platform Bring-Up \& Troubleshooting:} Led troubleshooting of critical platform bring-up faults, tracing panics to microcode defects and resolving PCIe device initialization errors linked to UEFI configuration settings.
\item \textbf{Build System Integration:} Maintain Yocto recipes and integrate software components into embedded Linux distributions.
\end{itemize}
\vspace{5pt}
\newpage
\textbf{Ericsson} \hfill Stockholm, Sweden\\
\job{HW \& SW Virtualization Developer}{July 2017 -- October 2021}{Pre-Silicon Hardware Emulation}{}
\begin{itemize}
\item \textbf{Pre-Silicon Simulation Design:} Architected high-quality simulation platforms for embedded Linux-based targets based on deep analysis of hardware specifications and hardware/software interfaces.
\item \textbf{Simulation IP Development:} Designed and implemented reusable simulation IPs including FPGA device models, interface controllers (PCIe, I2C, SPI), power/clock management controllers, microcontrollers, retimers, SFP modules, and GPIO controllers.
\item \textbf{Third-Party IP Integration:} Integrated third-party simulation IP, including PCIe controllers and instruction-set simulators (ISS), into Ericsson pre-silicon simulation environments.
\item \textbf{Standards Contribution:} Actively contributed to the evolution of the SystemC TLM standard through Accellera Systems Initiative working groups.
\end{itemize}
\vspace{5pt}
\job{Master Thesis Student}{January 2017 -- June 2017}{Thesis For M.S. In Embedded Systems}{}
\begin{itemize}
\item Researched and implemented a customized memory-error detection tool tailored for a proprietary DSP architecture, significantly improving runtime software stability.
\end{itemize}
\vspace{5pt}
\textbf{Imagimob} \hfill Stockholm, Sweden\\
\job{Firmware Developer}{April 2016 -- June 2016}{Embedded Firmware}{}
\begin{itemize}
\item Engineered low-latency multi-sensor data fusion firmware optimized for memory- and power-constrained edge computing hardware platforms.
\end{itemize}
\section{Professional Training}
\begin{itemize}
\item Linux Kernel Internals \& Development (LFD420) -- The Linux Foundation.
\item Developing Linux Device Drivers (LFD430) -- The Linux Foundation.
\item Security \& The Linux Kernel (LFD441) -- The Linux Foundation.
\item Wind River Linux LTS Application Development -- Wind River.
\item C \& C++ SEI CERT Based Secure Coding (ARM) -- SCADEMY Human-AI Symbiosis Academy
\end{itemize}
\section{Education}
\textbf{M.S. in Embedded Systems} \hfill August 2015 -- December 2017\\
KTH Royal Institute of Technology, Sweden\par\vspace{4pt}
\textbf{B.Tech. in Instrumentation \& Control} \hfill July 2011 -- June 2015\\
Manipal Institute of Technology, India
\section{Languages}
English \quad Hindi \quad Bengali \quad Swedish \quad Gujarati
\end{document}
